% =============================================================
%  RESULTS  (sections/04_results.tex)
% =============================================================

\chapter{Final Network Configuration and Production Strategy}

% Present your results using the SAME subsection structure as the
% Procedure chapter. Reference figures and tables in the text;
% do not let them stand alone without discussion.

% ----- Subsection 1 -------------------------------------------
\section{Wells}
\label{sec:res_wells}
The well design process was developed as shown in Section~\ref{subsec:well}. It was defined the tubing and production casing for the default well. 
It was also informed the gas lift rate needed to mantain the production rate of 6428.57 bbl/day. The final design of the well is shown in Table~\ref{tab:well_design}.
\begin{table}[ht]
    \centering
    \caption{Final well design.}
    \label{tab:well_design}
    \begin{tabular}{lccc}
        \hline
        Parameter & OD & ID & Unit \\
        \hline
        Tubing Diameter & 3.5 & 2.922 & inches \\
        Production Casing Diameter & 7.0 & 5.92 & inches \\
        \hline
    \end{tabular}
\end{table}

The gas lift is need for water cuts higher than 5\%. The maximum gas lift rate needed is 1.6~MMscf/day, for a water cut of 95\%. The Table~\ref{tab:gas_lift} shows the minimum gas lift rate needed for different water cuts.
Using those gas lift rates is also asured that the well head pressure will be higher than 1200 psi, which is the minimum pressure needed to avoid parafin deposition.
\begin{table}[ht]
    \centering
    \caption{Minimum gas lift rate needed for different water cuts.}
    \label{tab:gas_lift}
    \begin{tabular}{lcc}
        \hline
        Water Cut (\%) & Gas Lift Rate (MMscf/day) \\
        \hline
        2 & 0.0 \\
        5 & 0.0 \\
        20 & 0.2 \\
        35 & 0.6 \\
        70 & 1.2 \\
        90 & 1.6 \\ 
        95 & 1.6 \\
        \hline
    \end{tabular}
\end{table}

\section{Pipeline}

\section{Production Strategy}

It is expected that the field can produce with 3 wells off. If it is necessary to maintain the minimum well head pressure, the total 
field liquid production should be reduced to 19285 bbl/day, correponding to a 4 wells producing at the maximum liquid rate.
The economic analysis of the project shows that it more profitable to produce the well with a reduced gas lift injection rate. The Figure~\ref{optimal_gas_lift_50} and Figure~\ref{optimal_gas_lift_30} show the optimal gas lift injection rate for different 
oil prices, for a \$50/bbl and \$30/bbl, respectively. 

\section{Network Simulator}

A network simulator was developed to model the production of the field. It used tha Hagedorn and Brown correlation to model the
two-phase flow in the well and the Beggs and Brill correlation to model the two-phase flow in the pipeline. 
The simulator was used to predict the flow rates and pressures in the wells and manifolds, showing a good agreement with the 
Pipesim results, being a useful tool for the production optimization of the field in the case where Pipesim is not available.


